\section*{Chapter \ref{ch:b1:curves} --- Plane Curves}

\begin{solution}{exo:b1:curves:1}
$M(-t) = (t^2, -t^3)$: reflection in the $x$-axis; study $t \geq 0$.
Both $x' = 2t$ and $y' = 3t^2$ are $\geq 0$: the branch moves
right and up, from $(0,0)$ to infinity. At $t = 0$ the velocity
vanishes; expansions $x = t^2$, $y = t^3$ show the tangent is the
$x$-axis ($y/x = t \to 0$) with $y$ changing sign while $x \geq 0$:
a \emph{\omterm{ex:b1:curves:astroid}{cusp}} pointing left. The curve is the semi-cubical parabola
$y^2 = x^3$.
\end{solution}

\begin{solution}{exo:b1:curves:2}
$M(t + 2\pi) = M(t) + (2\pi, 0)$: the curve repeats, translated by
one wheel circumference; study one period. $x'(t) = 1 - \cos t \geq
0$, $y'(t) = \sin t$: the arch rises on $\intoo{0}{\pi}$, falls on
$\intoo{\pi}{2\pi}$, culminating at $(\pi, 2)$. Singular points
where $x' = y' = 0$: $t \in 2\pi\Z$, on the ground. Near $t = 0$:
\[
x(t) = \frac{t^3}{6} + o(t^3),
\qquad
y(t) = \frac{t^2}{2} + o(t^2):
\]
$x$ changes sign, $y \geq 0$, and $\frac{x}{y^{3/2}}$ bounded: the
tangent is \emph{vertical} (direction $(0,1)$), a \omterm{ex:b1:curves:astroid}{cusp} where the
tracked point momentarily has zero speed --- the physical signature
of rolling without slipping.
\end{solution}

\begin{solution}{exo:b1:curves:3}
At $t = \frac\pi4$: $M = \bigl(\frac{\sqrt2}{4},
\frac{\sqrt2}{4}\bigr)$, tangent direction $(-\cos t, \sin t) =
\frac{1}{\sqrt2}(-1, 1)$ (\cref{ex:b1:curves:astroid}). Tangent
line: $y - \frac{\sqrt2}{4} = -(x - \frac{\sqrt2}{4})$, i.e.\ $x +
y = \frac{\sqrt2}{2}$. Intercepts: $\bigl(\frac{\sqrt2}{2},
0\bigr)$ and $\bigl(0, \frac{\sqrt2}{2}\bigr)$; the segment between
them has \omterm{def:b1:curves:arclength}{length} $\sqrt{\frac12 + \frac12} = 1$.

(General point $t$: the tangent at $(\cos^3 t, \sin^3 t)$ cuts the
axes at $(\cos t, 0)$ and $(0, \sin t)$ --- check that the line
through these points has direction $(-\cos t, \sin t)$ and passes
through $M(t)$ --- and the cut segment has \omterm{def:b1:curves:arclength}{length} $\sqrt{\cos^2 t +
\sin^2 t} = 1$.)
\end{solution}

\begin{solution}{exo:b1:curves:4}
$r = \cos 2\theta$: period $\pi$ in $\theta$, symmetric in both
axes; $r$ vanishes at $\theta = \pm\frac\pi4$ (tangents at the
origin along the diagonals) and is negative for $\theta \in
\intoo{\frac\pi4}{\frac{3\pi}{4}}$, where the points plot on the
opposite ray --- producing four petals along the axes directions
$\theta = 0, \frac\pi2, \pi, \frac{3\pi}{2}$, each of maximal
radius $1$.

$r\cos\theta = 1$ is the equation $x = 1$: the \omterm{def:b1:curves:polar}{polar curve} $r =
\frac{1}{\cos\theta}$ is the vertical line $x = 1$ (swept once for
$\theta \in \intoo{-\frac\pi2}{\frac\pi2}$).
\end{solution}

\begin{solution}{exo:b1:curves:5}
$r(\theta) = 1 + 2\cos\theta$ vanishes for $\cos\theta = -\frac12$:
$\theta = \pm\frac{2\pi}{3}$. Symmetry in the $x$-axis; study
$\theta \in \intcc{0}{\pi}$. For $\theta \in
\intco{0}{\frac{2\pi}{3}}$: $r > 0$, decreasing from $3$ to $0$:
outer arc, entering the origin tangent to the ray $\theta =
\frac{2\pi}{3}$. For $\theta \in \intoc{\frac{2\pi}{3}}{\pi}$: $r <
0$: the points $r\vec u(\theta)$ sit on the opposite ray (angle
$\theta - \pi \in \intoc{-\frac\pi3}{0}$), with distance $\abs r$
growing from $0$ to $1$: this draws a small \emph{inner loop}
through the origin. At $\theta = \pi$, $r = -1$ and the point is
$-\vec u(\pi) = (1, 0)$: the loop closes on the positive $x$-axis.
Sketch: a big heart-like outer curve
with maximum reach $3$ at $\theta = 0$, plus a loop inside it
through $O$ and $(1,0)$.
\end{solution}

\begin{solution}{exo:b1:curves:6}
Complete squares:
\[
9(x^2 - 4x) + 25(y^2 - 2y) = 164
\iff 9(x-2)^2 + 25(y-1)^2 = 164 + 36 + 25 = 225 .
\]
Dividing by $225$: $\frac{(x-2)^2}{25} + \frac{(y-1)^2}{9} = 1$: an
ellipse of center $(2, 1)$, semi-axes $a = 5$ (horizontal), $b =
3$; $c = \sqrt{25 - 9} = 4$: foci $(2 \pm 4,\, 1) = (-2, 1)$ and
$(6, 1)$; eccentricity $e = \frac45$.
\end{solution}

\begin{solution}{exo:b1:curves:7}
Focus at the origin, directrix $D: x = d$ with $d = \frac pe > 0$.
For a point $M = (r\cos\theta, r\sin\theta)$ with $r > 0$:
\[
MF = r,
\qquad
d(M, D) = \abs{d - r\cos\theta},
\]
and the \omterm{def:b1:curves:conics}{conic} condition $MF = e\,d(M, D)$, in the regime where $M$
is on the focus side of the directrix ($r\cos\theta < d$), reads $r
= e(d - r\cos\theta)$, i.e.
\[
r(1 + e\cos\theta) = ed = p,
\qquad
r = \frac{p}{1 + e\cos\theta} .
\]
For $e < 1$ the denominator never vanishes: all $\theta$ allowed
(ellipse). For $e = 1$: $\theta \neq \pi$ (parabola, open towards
the directrix's far side). For $e > 1$: need $1 + e\cos\theta > 0$,
i.e.\ $\theta \in \intoo{-\theta_0}{\theta_0}$ with $\theta_0 =
\arccos\bigl(-\frac1e\bigr)$: one \emph{branch} of the hyperbola
(the other branch corresponds to the sign choice $r < 0$, or to the
second focus).
\end{solution}

\begin{solution}{exo:b1:curves:8}
\emph{Gardener:} attach a string of \omterm{def:b1:curves:arclength}{length} $2a$ to two stakes $F,
F'$ and keep it taut with the tracing point $M$: the constraint is
exactly $MF + MF' = 2a$, so the traced curve is the ellipse.

\emph{Reflection property:} let $t \mapsto M(t)$ be a regular
parametrization and $u(t) = \frac{M(t) - F}{\norm{M(t) - F}}$,
$u'(t)$ the analogous unit vector to $F'$. Differentiating
$\norm{M - F} + \norm{M - F'} = 2a$, using $\frac{\dd}{\dd
t}\norm{M - F} = \bigl\langle \frac{M - F}{\norm{M-F}},\,
M'\bigr\rangle$ (chain rule on $\sqrt{\langle\cdot,\cdot\rangle}$):
\[
\bigl\langle u + u',\, M'\bigr\rangle = 0 .
\]
So the tangent direction $M'$ is \omterm{def:b1:euclid:orthogonal}{orthogonal} to the bisector
direction $u + u'$ of the two focal rays: the tangent makes equal
angles with $MF$ and $MF'$. (A light ray from one focus reflects
off the ellipse to the other focus.)
\end{solution}

\begin{solution}{exo:b1:curves:9}
\begin{enumerate}
  \item Domain: $\cos 2\theta \geq 0$: $\theta \in
        \intcc{-\frac\pi4}{\frac\pi4} \cup
        \intcc{\frac{3\pi}{4}}{\frac{5\pi}{4}}$. Symmetries: $\theta
        \mapsto -\theta$ ($x$-axis) and $\theta \mapsto \pi -
        \theta$ ($y$-axis): study $\intcc{0}{\frac\pi4}$; $r$
        decreases from $1$ to $0$. At $\theta = \pm\frac\pi4$: $r =
        0$, tangents at the origin along the diagonals. The curve is
        the $\infty$ symbol: two symmetric loops meeting at $O$,
        reaching $(\pm 1, 0)$.
  \item With $F, F' = (\pm c, 0)$, $c = \frac{1}{\sqrt 2}$, and $M
        = (r\cos\theta, r\sin\theta)$:
        \[
        MF^2\, MF'^2
        = \bigl((r\cos\theta - c)^2 + r^2\sin^2\theta\bigr)
        \bigl((r\cos\theta + c)^2 + r^2\sin^2\theta\bigr)
        = (r^2 + c^2)^2 - (2rc\cos\theta)^2 ,
        \]
        by the identity $(A - B)(A + B) = A^2 - B^2$ with $A = r^2
        + c^2$, $B = 2rc\cos\theta$. With $c^2 = \frac12$:
        \[
        MF^2 MF'^2 = \Bigl(r^2 + \frac12\Bigr)^2 - 2r^2\cos^2\theta
        = r^4 + r^2\bigl(1 - 2\cos^2\theta\bigr) + \frac14
        = r^4 - r^2\cos 2\theta + \frac14 .
        \]
        On the lemniscate, $r^2 = \cos 2\theta$: the first two terms
        cancel, leaving $MF^2MF'^2 = \frac14$, i.e.\ $MF \cdot MF' =
        \frac12$. Conversely, the computation read backwards shows
        the locus equation $MF\,MF' = \frac12$ is $r^2 = \cos
        2\theta$ (for $r \neq 0$; and $O$ satisfies both).
\end{enumerate}
\end{solution}

\begin{solution}{exo:b1:curves:10}
$r = 1 + \cos\theta$, $r' = -\sin\theta$:
\[
r^2 + r'^2 = 1 + 2\cos\theta + \cos^2\theta + \sin^2\theta
= 2 + 2\cos\theta = 4\cos^2\frac\theta2 ,
\]
so $\sqrt{r^2 + r'^2} = 2\abs{\cos\frac\theta2}$ and, by the
symmetry in the $x$-axis,
\[
L = \int_0^{2\pi} 2\Bigl|\cos\frac\theta2\Bigr|\,\dd\theta
= 2\int_0^{\pi} 2\cos\frac\theta2\,\dd\theta
= 2\Bigl[4\sin\frac\theta2\Bigr]_0^{\pi} = 8 .
\]
Another algebraic, $\pi$-free perimeter.
\end{solution}

\begin{solution}{exo:b1:curves:11}
The point $(a\cos t, b\sin t)$ satisfies the equation:
$\cos^2 t + \sin^2 t = 1$. The line's normal vector is
$\bigl(\frac{\cos t}a, \frac{\sin t}b\bigr)$, and its product
with the velocity $(-a\sin t, b\cos t)$ is $-\sin t\cos t +
\sin t\cos t = 0$: the line passes through the point with the
tangent direction --- it is the tangent. For $(x_0, y_0)$ on the
ellipse, write $\cos t = \frac{x_0}a$, $\sin t = \frac{y_0}b$ and
substitute:
\[
\frac{x\,x_0}{a^2} + \frac{y\,y_0}{b^2} = 1 ,
\]
obtained from the ellipse equation by ``splitting'' $x^2 \mapsto
x\,x_0$ and $y^2 \mapsto y\,y_0$.
\end{solution}

\begin{solution}{exo:b1:curves:12}
\begin{enumerate}
  \item $r' = k\eu^{k\theta}$, so $\tan V = \frac{r}{r'} =
        \frac1k$: constant. The spiral cuts every ray from the
        origin at the same angle $V = \arctan\frac1k$.
  \item In complex notation the spiral is $\{\eu^{k\theta}
        \eu^{\iu\theta} : \theta \in \R\}$. Rotating by $c$
        multiplies by $\eu^{\iu c}$:
        \[
        \eu^{k\theta}\eu^{\iu(\theta + c)}
        = \eu^{-kc}\,\eu^{k(\theta + c)}\eu^{\iu(\theta+c)} ,
        \]
        and as $\theta + c$ runs over $\R$ this describes
        $\eu^{-kc}$ times the spiral: rotation $=$ scaling. No
        other smooth curve but lines and circles has this
        property.
  \item $\sqrt{r^2 + r'^2} = \sqrt{1 + k^2}\,\eu^{k\theta}$, so
        on $\intcc{A}{\theta_0}$ the \omterm{def:b1:curves:arclength}{length} is
        $\frac{\sqrt{1+k^2}}{k}\bigl(\eu^{k\theta_0} -
        \eu^{kA}\bigr)$, and as $A \to -\infty$:
        \[
        L = \frac{\sqrt{1 + k^2}}{k}\,\eu^{k\theta_0} < \infty :
        \]
        infinitely many turns around the origin, finite total
        \omterm{def:b1:curves:arclength}{length} (the turns shrink geometrically).
\end{enumerate}
\end{solution}

\begin{solution}{pb:b1:curves:1}
\textbf{1.} Rolling without slipping means the contact point has
traveled a distance equal to the arc of wheel unrolled: after
turning by $t$, the center is at $(Rt, R)$. The marked point sits
on the rim at angle $t$ \emph{behind} the downward vertical
(the wheel turns clockwise while advancing):
\[
M(t) = \Omega(t) + R(-\sin t, -\cos t)
= \bigl(R(t - \sin t),\ R(1 - \cos t)\bigr),
\]
which is correct at $t = 0$ ($M = (0,0)$) and at $t = \pi$ ($M =
(\pi R, 2R)$, the top).

\textbf{2.} $f'(t) = R(1 - \cos t,\ \sin t)$ and
\[
\norm{f'}^2 = R^2\bigl((1-\cos t)^2 + \sin^2 t\bigr)
= 2R^2(1 - \cos t) = 4R^2\sin^2\frac t2 ,
\]
so $\norm{f'} = 2R\abs{\sin\frac t2}$. Singular points at $t \in
2\pi\Z$: the \omterm{ex:b1:curves:astroid}{cusps} on the ground (\cref{exo:b1:curves:2}); the
top of the arch is $t = \pi$, where the speed $2R$ is maximal.

\textbf{3.} Half-angle factorizations:
\[
f'(t) = 2R\sin\frac t2\,\Bigl(\sin\frac t2,\ \cos\frac t2\Bigr),
\quad
M - C = 2R\sin\frac t2\,\Bigl(-\cos\frac t2,\ \sin\frac t2\Bigr),
\]
\[
T - M = \bigl(R\sin t,\ R(1 + \cos t)\bigr)
= 2R\cos\frac t2\,\Bigl(\sin\frac t2,\ \cos\frac t2\Bigr).
\]
For $0 < t < 2\pi$, $\sin\frac t2 \neq 0$: $M - C$ is \omterm{def:b1:euclid:orthogonal}{orthogonal}
to the tangent direction $\bigl(\sin\frac t2, \cos\frac
t2\bigr)$ (their product is $-\sin\frac t2\cos\frac t2 +
\sin\frac t2\cos\frac t2 = 0$), and $T - M$ is parallel to it.
The chord to the top of the wheel \emph{is} the tangent; the
chord to the contact point is the normal.

\textbf{4.} At each instant the wheel pivots about its contact
point (that point has zero velocity: rolling without slipping),
so every rigid point of the wheel moves \omterm{def:b1:euclid:orthogonal}{orthogonally} to the line
joining it to $C$, with speed (angular velocity $1$) equal to
that distance. Check: $\norm{M - C} = 2R\abs{\sin\frac t2} =
\norm{f'(t)}$.

\textbf{5.} $2R\,y(t) = 2R\cdot 2R\sin^2\frac t2 =
4R^2\sin^2\frac t2 = \norm{f'(t)}^2$. The speed at height $y$ is
$\sqrt{2R\,y}$ --- formally the law $v = \sqrt{2gh}$ of \omterm{def:b1:vspaces:free}{free}
fall, with the ground playing the ceiling; Part IV turns this
observation into clockwork.

\textbf{6.} By \cref{def:b1:curves:arclength} and question 2
($\sin\frac t2 \geq 0$ on $\intcc{0}{2\pi}$):
\[
L = \int_0^{2\pi} 2R\sin\frac t2\,\dd t
= 2R\Bigl[-2\cos\frac t2\Bigr]_0^{2\pi} = 2R(2 + 2) = 8R .
\]
Wren's theorem: four diameters exactly.

\textbf{7.} $s(t) = \int_0^t 2R\sin\frac u2\,\dd u =
4R\bigl(1 - \cos\frac t2\bigr)$; $s(2\pi) = 4R(1+1) = 8R$,
consistent.

\textbf{8.} $\sigma(t) = \abs{s(t) - s(\pi)} = \abs{4R(1 -
\cos\frac t2) - 4R} = 4R\abs{\cos\frac t2}$, so $\sigma^2 =
16R^2\cos^2\frac t2$; and $2R - y = R(1 + \cos t) =
2R\cos^2\frac t2$, whence $8R(2R - y) = 16R^2\cos^2\frac t2 =
\sigma^2$.

\textbf{9.} At the \omterm{ex:b1:curves:astroid}{cusp} $t = 0$ (or $2\pi$): $\sigma = 4R$, half
of $8R$: the apex halves the arch. In both computations the
speed is $\abs{\text{trigonometric polynomial in } t/2}$, whose
antiderivative is again trigonometric: the \omterm{def:b1:curves:arclength}{length} is a
difference of \emph{values} of cosines --- rational numbers
times $R$ --- with no arc of circle to measure, hence no $\pi$.
The circle itself has constant speed, so its \omterm{def:b1:curves:arclength}{length} \omterm{thm:b1:integration:def}{integral}
produces the full \omterm{prop:b1:reals:intervals}{interval} \omterm{def:b1:curves:arclength}{length} $2\pi$; the cycloid's speed
vanishes at the ends and integrates algebraically.

\textbf{10.} The arch is swept with $x$ increasing from $0$ to
$2\pi R$ ($x' = R(1 - \cos t) \geq 0$, vanishing only at
isolated points). Substituting $x = x(t)$ in the area \omterm{thm:b1:integration:def}{integral}
$\int_0^{2\pi R} y\,\dd x$ (\cref{thm:b1:integration:parts})
gives $A = \int_0^{2\pi} y(t)\,x'(t)\,\dd t$.

\textbf{11.}
\[
A = \int_0^{2\pi} R(1 - \cos t)\cdot R(1 - \cos t)\,\dd t
= R^2\int_0^{2\pi}\bigl(1 - 2\cos t + \cos^2 t\bigr)\dd t
= R^2\Bigl(2\pi - 0 + \pi\Bigr) = 3\pi R^2 ,
\]
using $\int_0^{2\pi}\cos^2 = \pi$. Exactly three wheel areas:
Galileo's balance said ``about $3$''; Roberval's computation
says ``exactly''.

\textbf{12.} With $x = \cos^3 t$, $y = \sin^3 t$: $y\,x' =
-3\sin^4 t\cos^2 t$. Linearize:
\[
\sin^4 t\cos^2 t = \frac{1 - \cos 2t}{2}\cdot\frac{\sin^2
2t}{4} = \frac{\sin^2 2t}{8} - \frac{\sin^2 2t\cos 2t}{8} ,
\]
and over $\intcc{0}{2\pi}$: $\int\sin^2 2t = \pi$, $\int \sin^2
2t\cos 2t = \bigl[\frac{\sin^3 2t}{6}\bigr] = 0$. So $\oint
y\,\dd x = -3\cdot\frac\pi8$, and the enclosed area is
$\frac{3\pi}8$ (the sign records the counterclockwise
orientation).

\textbf{13.} For $(\cos t, \sin t)$ with $t$ from $\pi$ to $0$,
$x$ increases from $-1$ to $1$ and
\[
\int_\pi^0 \sin t\cdot(-\sin t)\,\dd t
= \int_0^\pi \sin^2 t\,\dd t = \frac\pi2 :
\]
the formula returns the area of the upper half-disc, as it must.

\textbf{14.} $x' = R(1 + \cos t) = 2R\cos^2\frac t2$, $y' =
R\sin t = 2R\sin\frac t2\cos\frac t2$, so $\norm{f'} =
2R\cos\frac t2$ (nonnegative on $\intcc{-\pi}{\pi}$). Arc from
the bottom: $s(t) = \int_0^t 2R\cos\frac u2\,\dd u =
4R\sin\frac t2$. Then
\[
y = R(1 - \cos t) = 2R\sin^2\frac t2
= 2R\Bigl(\frac{s}{4R}\Bigr)^{\!2} = \frac{s^2}{8R} .
\]

\textbf{15.} Energy conservation with $v = \frac{\dd s}{\dd\tau}$
and $y = \frac{s^2}{8R}$, $y_0 = \frac{s_0^2}{8R}$:
\[
\Bigl(\frac{\dd s}{\dd\tau}\Bigr)^{\!2} = 2g(y_0 - y)
= \frac{2g}{8R}\bigl(s_0^2 - s^2\bigr)
= \frac{g}{4R}\bigl(s_0^2 - s^2\bigr).
\]

\textbf{16.} Differentiating in $\tau$: $2s's'' =
-\frac{g}{4R}\,2ss'$, so wherever $s' \neq 0$ (hence everywhere
by \omterm{def:b1:continuity:continuous}{continuity}), $s'' = -\frac{g}{4R}s$: the harmonic oscillator.
By \cref{thm:b1:diffeq:homogeneous2}, $s(\tau) =
A\cos\omega\tau + B\sin\omega\tau$ with $\omega =
\sqrt{g/(4R)}$; the initial conditions $s(0) = s_0$, $s'(0) = 0$
give $s(\tau) = s_0\cos\omega\tau$.

\textbf{17.} The bead reaches the bottom when $s = 0$, i.e.\ at
$\omega\tau = \frac\pi2$:
\[
T_{\downarrow} = \frac{\pi}{2\omega}
= \frac\pi2\sqrt{\frac{4R}{g}} = \pi\sqrt{\frac Rg}\,,
\]
independent of $s_0$: released from anywhere in the bowl, all
beads arrive at the same instant --- the tautochrone.

\textbf{18.} Substituting $s = s_0\cos\omega\tau$: the left side
is $s_0^2\omega^2\sin^2\omega\tau$ and the right side is
$\frac{g}{4R}s_0^2(1 - \cos^2\omega\tau) =
s_0^2\omega^2\sin^2\omega\tau$: exact equality, so no spurious
solution was introduced. For a circular arc, $y$ is not
proportional to $s^2$ ($y = \ell(1 - \cos\frac s\ell) =
\frac{s^2}{2\ell} - \frac{s^4}{24\ell^3} + \dots$): the
restoring term is only \emph{approximately} \omterm{def:b1:linmaps:def}{linear}, so the
period of a circular pendulum drifts with amplitude, while the
cycloid's is rigorously constant.

\textbf{19.} $\pi\sqrt{R/g} = 1$ gives $R = \frac{g}{\pi^2}
\approx \frac{9.81}{9.87} \approx 0.994$ m --- almost exactly
one meter. The full oscillation takes $\frac{2\pi}\omega =
2\pi\sqrt{4R/g}$, which is precisely the small-angle formula
$2\pi\sqrt{\ell/g}$ for a pendulum of \omterm{def:b1:curves:arclength}{length} $\ell = 4R$:
Huygens suspended his pendulum from cycloidal cheeks of exactly
that proportion.

\textbf{20.} At $t = \frac\pi2$, $R = 1$: $M = \bigl(\frac\pi2 -
1,\ 1\bigr)$, $T = \bigl(\frac\pi2,\ 2\bigr)$, so $T - M = (1,
1)$. And $f'(\frac\pi2) = (1 - 0,\ 1) = (1, 1)$: the chord to
the top is the velocity, as promised.

\textbf{21.} For the trochoid, $x'(t) = R - d\cos t$ and $y'(t)
= d\sin t$. If $d < R$: $x' \geq R - d > 0$, the point always
advances; the velocity never vanishes (its first component is
positive): no singular point, a regular wave. If $d > R$: $x'(0)
= R - d < 0 < R + d = x'(\pi)$, so the point moves backwards
near the contact instants and forwards elsewhere: the curve
crosses itself in loops. A point on the flange of a railway
wheel (below the rail head, $d > R$) travels backwards at every
turn.

\textbf{22.} The chord from the \omterm{ex:b1:curves:astroid}{cusp} $(\pi R, 2R)$ to the origin
has \omterm{def:b1:curves:arclength}{length} $L = R\sqrt{\pi^2 + 4}$ and slope angle $\alpha$ with
$\sin\alpha = \frac{2R}{L}$. Sliding from rest with constant
acceleration $g\sin\alpha$: $L = \frac12 g\sin\alpha\,T^2$, so
\[
T = \sqrt{\frac{2L}{g\sin\alpha}} = \sqrt{\frac{2L^2}{2Rg}}
= \frac{L}{\sqrt{Rg}} = \sqrt{\pi^2 + 4}\,\sqrt{\frac Rg}
\approx 3.72\sqrt{\frac Rg}\,,
\]
against $\pi\sqrt{R/g} \approx 3.14\sqrt{R/g}$ for the cycloid:
the curved path is faster. It is in fact the fastest possible
--- the brachistochrone --- a theorem of the calculus of
variations, beyond this volume.

\textbf{23.} $\dfrac{\dd y}{\dd x} = \dfrac{y'}{x'} =
\dfrac{\sin t}{1 - \cos t} = \cot\dfrac t2$ (half-angle
formulas). Then
\[
\frac{\dd^2y}{\dd x^2}
= \frac{\dd}{\dd t}\Bigl(\cot\frac t2\Bigr)\cdot\frac{1}{x'(t)}
= -\frac{1}{2\sin^2\frac t2}\cdot\frac{1}{2R\sin^2\frac t2}
= -\frac{1}{4R\sin^4\frac t2} < 0 :
\]
concave throughout the arch; as $t \to 0^+$ or $2\pi^-$ the
slope $\cot\frac t2 \to \pm\infty$: vertical tangents at the
\omterm{ex:b1:curves:astroid}{cusps}.

\textbf{24.} The chord direction is $T - M \parallel
\bigl(\sin\frac t2, \cos\frac t2\bigr)$, which tends to $(0, 1)$
as $t \to 0^{+}$: the tangent at the \omterm{ex:b1:curves:astroid}{cusp} is vertical ---
recovered from pure geometry, with no Taylor expansion.

\textbf{25.} (i) Wren's theorem, $L = 8R$; Roberval's area, $A =
3\pi R^2$ (Galileo's conjectured ratio $3$); Huygens'
tautochrone, $T_\downarrow = \pi\sqrt{R/g}$. (ii) Every
computation ran on the half-angle factorizations $f' =
2R\sin\frac t2(\sin\frac t2, \cos\frac t2)$ and their bowl
analogue --- one identity powering tangent, \omterm{def:b1:curves:arclength}{length} and clock
alike. (iii) The intrinsic relation $y = s^2/8R$ converted the
energy equation into $s'' = -\frac{g}{4R}s$, a \omterm{def:b1:linmaps:def}{linear} equation
with constant coefficients whose solutions are exactly
isochronous. (iv) Part I used the differential calculus of
\cref{ch:b1:derivative}, Part II--III the \omterm{thm:b1:integration:def}{integral} of
\cref{ch:b1:integration}, Part IV the differential equations of
\cref{ch:b1:diffeq} --- the cycloid is this book's curriculum
rolled into one curve.
\end{solution}
